\documentclass[11pt,a4paper]{article}

\usepackage[T1]{fontenc}
\usepackage[utf8]{inputenc}
\usepackage[polish]{babel}
\usepackage{lmodern}
\usepackage[margin=2.5cm]{geometry}
\usepackage{amsmath,amssymb,amsthm}
\usepackage{hyperref}

\title{Synchroniczny wybór lidera na cyklu przy przybliżonym $n$}
\author{Kacper Paciorek}
\date{\today}

\newtheorem{theorem}{Twierdzenie}

\begin{document}

\maketitle

\section{Postawienie problemu}
Problem omawiany w pracy to wybór lider w siecie o następujących założeniach:
\begin{enumerate}
    \item Topologia - sieć to skierowany cykl na $n$ wierzchołkach
    \item Obliczenia są synchroniczne. W jednej rundzie proces:
        \begin{enumerate}
            \item Odbiara wiadomość jeśli jest taka na wejściu
            \item Wykonuje obliczenia zmieniając swój stan
            \item Może wysłać jedną wiadomość do następnego procesu na cyklu, która będzie odczytana w następnej rundzie
        \end{enumerate}
    \item Algorytm nie zna dokładnego $n$. Procesy dostają jedynie $N$ (wszystkie procesy to samo), które spełnia $N \leq n \leq 2N-1$
    \item Aby przełamać symetrię procesy mają dostęp do niezależnych generatorów liczb losowych
\end{enumerate}

Można się zastanawiać dlaczego zakładamy że znamy $n$ z taką dokładnością a nie inną. Odpowiedź na to pytanie wynika z poniższego.

\begin{theorem}
    Nie istnieje algorytm poprawnie rozwiązujący ten problem dla cyklu wielkości $N$ oraz $2N$.
\end{theorem}
\begin{proof}
    Załóżmy że mamy algorytm, który działa dla $N$. Rozważmy konkretny jego przebieg aż do zatrzymania. Przyjmijmy teraz, że 
    dla dwukrotnie dłuższego cyklu procesy $i, i+n$ przebiagają dokładnie tak samo jak przebiegało $i$ w przebiegu dla $N$. Jest to możliwe bo 
    generatory losowe są niezależne a topologia sieci sprawia w takiej sytuacji że procesy dostają te same wiadomości. To jednak oznacza, że na koniec 
    dwa procesy mają status lidera. Jest to sprzeczność.
\end{proof}

Z powyższego faktu widać, że założenie o aproksymacji $n$ jest najsłabsze przy którym taki algortym może istnieć. 

Dodatkowym problemem, który rozważają autorzy w tej sytuacji jest ustalenie dokładnej wartości $n$.

\section{Algorytm i poprawność}

Komunikacja w przedstawionym przez autorów algorytmie jest bardzo prosta. Sprowadza się ona tylko do wysyłania jednobitowej wiadomości. 
Procesy patrzą tylko na to czy w danej turze otrzymały wiadomość. 

Algorytm przedstawiony przez autorów zależy od parametru $c > 0$ i składa się z następujących faz:
\begin{enumerate}
    \item \textbf{Losowe wygenerowanie kandydatów na lidera.}
        \begin{itemize}
            \item Każdy z procesów z pradwopodobieństwem $\frac{c}{N}$ zostaje kandydatem na lidera
            \item Następnie procesy przez $2N-1$ rund komunikują się ze sobą w następujący sposób:
                \begin{itemize}
                    \item Kandydaci - inicjują komunikację ale nie przekazują danej otrzymanych wiadomości
                    \item Pozostałe procesy - jeśli w danej turze otrzymały wiadomość to wysyłają ją w następnej
                \end{itemize}
                Skoro $2N-1 \geq 1$ to po zakończeniu tej komunikacji proces zobaczył wiadomość dokładnie wtedy gdy został wylosowany jakiś kandydat.
                Zatem po zakończeniu tej fazy wszystkie procesy wiedzą czy istnieje kandydat.
            \item Jeśli nie został wybrany żaden kandydat procedura jest powtarzana aż do skutku
        \end{itemize}
    \item \textbf{Sprawdzenie czy został wylosowany dokładnie jeden kandydat.} składa się ono z dwóch podfaz: 
        \begin{enumerate}
            \item[2.1] \textbf{Co najmniej jeden proces ma się zorientować jeśli jest wiecej niż jeden kandydat.}
                W tym celu przez dokładnie $N-1$ rund procesy komunikują się zgodnie z regułą:
                \begin{itemize}
                    \item Kandydaci - inicjują komunikację ale nie przekazują dalej wiadomości
                    \item Pozostałe procesy - jeśli w danej turze otrzymały wiadomość to wysyłają ją w następnej
                \end{itemize}
                Jeśli, któryś z kandydatów otrzymał wiadomość to wie, że istnieje inny kandydat (n > N-1 więc nie mógł otrzymać spowrotem wiadomości wysłanej przez siebie). 
                Z drugiej strony jeśli jest więcej niż jeden kandydat to z zasady szufladkowej musi być dwóch w odległości co najwyżej $N-1$. To oznacza, że po zakończeniu tej fazy wiemy, że jest więcej niż
                jeden kandydat dokładnie wtedy gdy istnieje proces który o tym wie
            \item[2.2] \textbf {Informacja o tym czy kandydaci się powtórzyli zostaje rozpowszechniona.} W tym celu przez $2N-1$ rund procesy komunikują się tak samo jak w poprzedniaj fazie 
            z tą róźnićą, że komunikacja jest inicjowana tylko przez te procesy, które zorientowały się, że jest więcej niż jeden kandydat. Jako, że $2N-1 \geq n$ to po zakończeniu tej fazy ta informacja zdążyła dotrzeć do wszystkich procesów.
        \end{enumerate}
        Jeśli jest dokładnie jeden kandydat to procedura zakończyła się sukcesem. W przeciwnym wypadku powtarzamy cały algorytm aż do skutku.
\end{enumerate}

Z powyższego opisu algorytmu jasno wynika, że jeśli procedura ta się zakończyła to został wybrany dokładnie jeden lider. 

\section{Analiza złożoności}

Analizujemy dwie miary: \emph{złożoność bitową} $\bar{C}$, czyli oczekiwaną łączną liczbę bitów wysłanych przez wszystkie procesy,
oraz \emph{złożoność czasową} $\bar{T}$, czyli oczekiwaną liczbę rund do zakończenia algorytmu. Ponieważ każda wiadomość to pojedynczy bit,
złożoność bitowa jest równa łącznej liczbie wysłanych wiadomości. Wyniki zależą od parametru $c$ oraz od $N$ i $n$, przy czym $c < N$ i $N \leq n \leq 2N-1$.

Wprowadźmy oznaczenia:
\begin{itemize}
    \item \emph{Wywołanie} algorytmu to jedno przejście przez fazy 1, 2.1 i 2.2. Algorytm jest wywoływany od nowa, dopóki wywołanie nie zakończy się sukcesem.
    \item $\bar{k}_{SEL}$ - oczekiwana liczba iteracji fazy 1 w jednym wywołaniu (czyli aż wylosowany zostanie choć jeden kandydat),
    \item $\bar{k}_{ALG}$ - oczekiwana liczba wywołań algorytmu aż do sukcesu.
\end{itemize}

\subsection{Rozbicie na fazy}

Kluczową obserwacją jest to, że w każdej fazie liczba wysłanych bitów jest (prawie) deterministyczna,
a losowość pojawia się tylko w liczbie powtórzeń.

\paragraph{Faza 1.} Iteracja, w której nie wylosowano żadnego kandydata, nie generuje żadnej komunikacji. Gdy istnieje co najmniej jeden kandydat,
to bit wysłany przez kandydata biegnie aż do następnego kandydata (który go nie przekazuje) lub z powrotem do niego samego. Ponieważ $2N-1 \geq n$ rund wystarcza,
aby tak się stało, każda krawędź cyklu zostaje użyta dokładnie raz. Zatem niezależnie od liczby iteracji faza 1 kosztuje łącznie dokładnie $n$ bitów.
Czasowo jedna iteracja trwa $2N-1$ rund, więc faza 1 trwa łącznie $k_{SEL}(2N-1)$ rund, gdzie $k_{SEL}$ jest (losową) liczbą iteracji.

\paragraph{Faza 2.1.} Trwa dokładnie $N-1$ rund.
\begin{itemize}
    \item Jeśli jest jeden kandydat, to jego bit pokonuje $N-1 < n$ krawędzi i nie wraca do niego. Koszt: $N-1$ bitów.
    \item Jeśli kandydatów jest więcej, to każdy bit zatrzymuje się najpóźniej na następnym kandydacie, więc żadna krawędź nie jest użyta więcej niż raz. Koszt: co najwyżej $n$ bitów.
\end{itemize}

\paragraph{Faza 2.2.} Trwa dokładnie $2N-1$ rund.
\begin{itemize}
    \item Jeśli jest jeden kandydat, to nikt nie inicjuje komunikacji. Koszt: $0$ bitów.
    \item Jeśli kandydatów jest więcej, to istnieje co najmniej jeden proces inicjujący i (analogicznie jak w fazie 1, bo $2N-1 \geq n$) każda krawędź zostaje użyta dokładnie raz. Koszt: $n$ bitów.
\end{itemize}

Podsumowanie dla jednego wywołania:
\begin{center}
    \renewcommand{\arraystretch}{1.3}
    \begin{tabular}{lccc}
        \hline
        Faza & jeden kandydat & więcej niż jeden kandydat & liczba rund \\
        \hline
        1 (wszystkie iteracje) & $n$         & $n$        & $k_{SEL}(2N-1)$ \\
        2.1                    & $N-1$       & $\leq n$   & $N-1$ \\
        2.2                    & $0$         & $n$        & $2N-1$ \\
        \hline
        Razem                  & $n + N - 1$ & $\leq 3n$  & $3N-2+k_{SEL}(2N-1)$ \\
        \hline
    \end{tabular}
\end{center}

\subsection{Wzory ogólne}

Kolejne wywołania algorytmu są niezależne i mają to samo prawdopodobieństwo sukcesu, więc liczba wywołań ma rozkład geometryczny.
Wszystkie wywołania poza ostatnim są nieudane i kosztują co najwyżej $3n$ bitów, a ostatnie kosztuje $n+N-1$. Z liniowości wartości oczekiwanej otrzymujemy
\begin{equation}\label{eq:C}
    \bar{C} \leq 3n(\bar{k}_{ALG} - 1) + n + N - 1 = (3\bar{k}_{ALG} - 2)n + N - 1.
\end{equation}
Każde wywołanie trwa $3N-2+k_{SEL}(2N-1)$ rund, a liczba iteracji fazy 1 ma w każdym wywołaniu ten sam rozkład, zatem
\begin{equation}\label{eq:T}
    \bar{T} = \bar{k}_{ALG}\bigl(3N-2+\bar{k}_{SEL}(2N-1)\bigr).
\end{equation}
Pozostaje oszacować $\bar{k}_{SEL}$ i $\bar{k}_{ALG}$.

\subsection{Oszacowanie $\bar{k}_{SEL}$}

Niech $M$ oznacza liczbę kandydatów wylosowanych w jednej iteracji fazy 1. Każdy z $n$ procesów jest kandydatem z prawdopodobieństwem $c/N$, niezależnie od pozostałych, więc
\[
    \Pr(M = i) = \binom{n}{i}\left(\frac{c}{N}\right)^{i}\left(1-\frac{c}{N}\right)^{n-i}.
\]
Iteracja kończy fazę 1, gdy $M > 0$. Ponieważ $n \geq N$, a $1 - c/N < 1$, oraz $1 - x < e^{-x}$, mamy
\[
    \Pr\nolimits_{SEL} = \Pr(M>0) = 1 - \left(1-\frac{c}{N}\right)^{n} \geq 1 - \left(1-\frac{c}{N}\right)^{N} > 1 - e^{-c}.
\]
Liczba iteracji jest zmienną geometryczną z prawdopodobieństwem sukcesu $\Pr_{SEL}$, więc
\begin{equation}\label{eq:ksel}
    \bar{k}_{SEL} = \frac{1}{\Pr_{SEL}} < \frac{1}{1-e^{-c}}.
\end{equation}

\subsection{Oszacowanie $\bar{k}_{ALG}$}

Wywołanie kończy się sukcesem dokładnie wtedy, gdy faza 1 zakończyła się z dokładnie jednym kandydatem, czyli gdy $M = 1$ pod warunkiem $M > 0$. Oznaczmy $q = 1 - c/N$. Wtedy
\[
    \Pr\nolimits_{ALG} = \Pr(M=1 \mid M>0) = \frac{\Pr(M=1)}{\Pr(M>0)}
    = \frac{n\frac{c}{N}\,q^{n-1}}{1-q^{n}}
    = \frac{n\frac{c}{N}}{(q^{-n}-1)\,q}.
\]
Stąd, korzystając kolejno z $\frac{N}{n}q < 1$ (bo $n \geq N$ i $q<1$), z $n \leq 2N-1$ oraz z $q^{-(2N-1)} < q^{-2N}$ (bo $q<1$):
\begin{equation}\label{eq:kalg1}
    \bar{k}_{ALG} = \frac{1}{\Pr_{ALG}} = \frac{N}{n}\,q\cdot\frac{q^{-n}-1}{c}
    < \frac{q^{-n}-1}{c} \leq \frac{q^{-(2N-1)}-1}{c} < \frac{\left(q^{-N}\right)^{2}-1}{c}.
\end{equation}

Oszacujmy teraz $q^{-N}$. Dla $c < N$, korzystając z rozwinięcia $-\ln(1-x) = \sum_{k \geq 1} x^k/k$ dla $x = c/N \in (0,1)$:
\begin{align*}
    \left(1-\frac{c}{N}\right)^{-N}
    &= \exp\left(-N\ln\left(1-\frac{c}{N}\right)\right) = \exp\left(N\sum_{k=1}^{\infty}\frac{1}{k}\left(\frac{c}{N}\right)^{k}\right) \\
    &= e^{c}\exp\left(N\sum_{k=2}^{\infty}\frac{1}{k}\left(\frac{c}{N}\right)^{k}\right)
    \leq e^{c}\exp\left(\frac{N}{2}\sum_{k=2}^{\infty}\left(\frac{c}{N}\right)^{k}\right) \\
    &= e^{c}\exp\left(\frac{N}{2}\cdot\frac{(c/N)^{2}}{1-c/N}\right)
    = e^{c}\exp\left(\frac{c^{2}}{2(N-c)}\right).
\end{align*}
Wstawiając to do \eqref{eq:kalg1} dostajemy
\begin{equation}\label{eq:kalg}
    \bar{k}_{ALG} < \frac{1}{c}\left[e^{2c}\exp\left(\frac{c^{2}}{N-c}\right) - 1\right].
\end{equation}

\subsection{Wynik}

Łącząc \eqref{eq:C}, \eqref{eq:T}, \eqref{eq:ksel} i \eqref{eq:kalg} otrzymujemy następujące twierdzenie.

\begin{theorem}
    Dla $0 < c < N$ algorytm wybiera lidera ze złożonością bitową i czasową
    \[
        \bar{C} \leq (3\bar{k}_{ALG}-2)n + N - 1, \qquad
        \bar{T} = \bar{k}_{ALG}\bigl(3N-2+\bar{k}_{SEL}(2N-1)\bigr),
    \]
    gdzie
    \[
        \bar{k}_{SEL} < \frac{1}{1-e^{-c}}, \qquad
        \bar{k}_{ALG} < \frac{1}{c}\left[e^{2c}\exp\left(\frac{c^{2}}{N-c}\right) - 1\right].
    \]
\end{theorem}

\paragraph{Zależność od $c$.} Parametr $c$ pozwala wymieniać czas na komunikację. Gdy $c \to 0$, to $\bar{k}_{ALG}$ zbiega do $2$, więc $\bar{C}$ zbiega do $n+N-1$,
ale jednocześnie $\bar{k}_{SEL} \to \infty$, a razem z nim $\bar{T}$. Z kolei duże $c$ powoduje, że $\bar{k}_{ALG}$ rośnie, bo z dużym prawdopodobieństwem wylosowanych zostaje wielu kandydatów.

\paragraph{Przypadek $c = \frac{1}{2}$.} Dla $N=1$ problem jest trywialny, więc przyjmijmy $N \geq 2$. Ograniczenie \eqref{eq:kalg} maleje wraz z $N$, więc wystarczy je policzyć dla $N=2$:
\begin{align*}
    \bar{k}_{ALG} &< 2\left(e\cdot\exp\left(\frac{1/4}{2-1/2}\right)-1\right) = 2\left(e^{7/6}-1\right) \approx 4{,}4225, \\
    \bar{k}_{SEL} &< \frac{1}{1-e^{-1/2}} \approx 2{,}5415.
\end{align*}
Zatem
\begin{align*}
    \bar{C} &\leq 3n(\bar{k}_{ALG}-1) + n + N - 1 < 11{,}28\,n + N - 1, \\
    \bar{T} &< 4{,}4225\cdot\bigl(3N-2+2{,}5415\,(2N-1)\bigr) < 35{,}75\,N - 20{,}08.
\end{align*}
W szczególności $\bar{C} < 12n + N - 1 = O(n)$ oraz $\bar{T} < 36N = O(n)$. Zarówno liczba wysłanych bitów, jak i liczba rund, są więc liniowe względem rozmiaru cyklu,
mimo że algorytm zna tylko przybliżenie $N$ liczby $n$.

Autorzy w pracy podają, że wyniki w tej sytuacji to $\bar{C} \leq 10n + N - 1$ oraz $\bar{T} \leq 35N$ jednak w dowodzie tego faktu występują błędy rachunkowe. 
Autorzy najpierw źle podstawiają jedną drugą do ograniczenia na $\bar{k}_{ALG}$ otrzymująć $4{,}2714295$ zamiast $4{,}4225$. 
Następnie w szacowaniu na $\bar{C}$ zamiast policzyć współczynnik przy $n$ jako $3\bar{k}_{ALG} + 1$ ignorują $+1$ dostając zbyt mało.
Błąd przy $\bar{T}$ jest konsekwencją zaniżonej wartości $\bar{k}_{ALG}$ .

\section{Symulacje}

\section{Dodatek - wyznaczenie $n$}
Warto jeszcze dodać, że po wyznaczeniu lidera można bardzo łatwo poznać dokładną wartość $n$. W tym celu lider wysyła wiadomość. 
Pozostałe procesory przekazują ją w kółko. Pozwalamy aby wiadomość obiegła cykl dwa razy. Wtedy każdy proces może obliczyć $n$ jako róźnica w numerze rundy
pomiędzy otrzymaniem wiadomości po raz pierwszy i po raz drugi.

\end{document}
